% Source-derived chapter 07: The main theorems, and a conjecture
% Original source: universal-p-adic-gross-zagier-formula
\chapter{The main theorems, and a conjecture}
\label{universal-p-adic-gross-zagier-formula-chapter-07}
\source{universal-p-adic-gross-zagier-formula}

\begin{remark}[Source section]
\label{universal-p-adic-gross-zagier-formula-chapter-07-source}
\source{universal-p-adic-gross-zagier-formula}
\source{universal-p-adic-gross-zagier-formula}
This chapter reproduces the corresponding section of the retrieved
LaTeX source, including its definitions, statements, examples, and
proofs. It has not yet been translated into Lean.
\notready
\end{remark}

% The source section title was: The main theorems, and a conjecture
In this section, we prove our main theorems (\S~\ref{sec: final}, or \S~\ref{sec G} for  Theorem \ref{nv exc}), as well as a universal Waldspurger formula for  families of `sign $+1$' (\S~\ref{sec: wald}). Then, we discuss a conjecture on the leading terms of universal Heegner points (and toric periods) at classical points (\S~\ref{sec bd conj}).






















 
 
 
 















   



\subsection{Proofs of the main theorems}\lb{sec: final}
Both of our central  theorems (Theorems \ref{GZ thm} and \ref{ugz thm}) ultimately follow from \cite{dd-pyzz, nonsplit}, where Theorem \ref{GZ thm} is established when $W$ is trivial, by an argument combining interpolation and multiplicity-one principles.


\subsubsection{$p$-adic Gross--Zagier formula for ordinary forms} \lb{B ord}
We start by stating a  variant of  Theorem \ref{GZ thm}, valid under  the same assumptions. 

  \begin{customthm}
\source{universal-p-adic-gross-zagier-formula}
\notready{$\text{B}^{\, \ord}$}\label{GZ thm'} Let $\Pi=\pi\ot \chi$ be an 
\source{universal-p-adic-gross-zagier-formula}
ordinary, locally distinguished automorphic  
 representation of $(\G\ts\H)'(\A)$ over $L$.
Let $V=V_{\Pi}$, and let $\wtil{P}_{\Pi}^{\ord} \in \wtil{H}^{1}_{f}(E, V_{\Pi})$  be the enhanced Heegner class defined in \eqref{tilPo}.

 Then 
for all $f_{1}^{}\in \Pi_{H'_{\infty}}^{\ord}$, $f_{2}^{}\in\Pi^{\vee, \ord}_{H'_{{\infty}}}$, $f_{3}^{}\in \Pi^{\ord}$, $f_{4}^{}\in\Pi^{\vee, \ord}$ with $(f_{3}, f_{4})^{\ord}\neq 0$, we have
\beq\lb{formula B'}   {  h_{V}(\wtil P_{\Pi}^{\ord}(f_{1}^{}),\wtil P_{ \Pi^{\vee}}^{\ord}(f_{2}^{}))    \over   (f_{3}^{}, f_{4}^{})^{\ord}_{\Pi} }  
= 
   \calL_{p}'(V_{(\pi, \chi)}, 0)
\cdot
Q^{\ord}\left( {  f_{1} \ot  f_{2}\over f_{3}\ot f_{4}}\right).
\eeq
\end{customthm}
\begin{rema}
\source{universal-p-adic-gross-zagier-formula}
\notready  In contrast to Theorem \ref{GZ thm}:
\begin{itemize}
\item
Theorem \ref{GZ thm'} also holds for exceptional $\Pi$;
\item we have only included the Gross--Zagier formula and omitted an analogue to the first statement of Theorem \ref{GZ thm}, that is that $\wtil P_{\Pi}^{\ord}$ takes values in $\wtil H^{1}_{f}(E, V)$, as that has already been established. 
\end{itemize} 
\end{rema}



\begin{lemm} \lb{B Bord} Suppose that $\Pi$ is not exceptional. Theorem \ref{GZ thm'} is equivalent to Theorem \ref{GZ thm}.
\end{lemm}
\begin{proof}  Using freely the notation and results of Appendix  \ref{app A}, we first show that Theorem \ref{GZ thm'} for $f_{1}$, $f_{2}$, $f_{3}$, $f_{4}$ is equivalent to  Theorem B for 
$$f_{1}'= \gamma_{H'}^{\ord} f_{1} \in \Pi_{H'_{p\infty}}, \quad f_{2}'= \gamma_{H'}^{\ord} f_{2}^{} \in \Pi^{\vee}_{H'_{p\infty}},\qquad f_{3}'=w_{\rm a}^{\ord} f_{3}^{}\in \Pi^{\rm a}, \quad f_{4}'=f_{4}^{}\in \Pi^{\vee, \ord};$$
let us call such $(f_{1}', f_{2}', f_{3}', f_{4}' )$  a `special quadruple'.

Indeed, by the definitions \eqref{Pord text}, \eqref{()'}, the left hand side of \eqref{formula B'} equals 
$$  {  h_{V}(P_{\Pi}^{}(f_{1}'), P_{\Pi^{\vee}}^{\ord}(f_{2}')) \over \dim W \cdot (f_{3}', f_{4}')};$$
whereas by  Proposition \ref{compare Qs}, 
\beqq
Q^{\ord}\left( {  f_{1} \ot  f_{2}\over f_{3}\ot f_{4}}\right).
= {e_{p}(V_{(\pi, \chi)})^{-1} \over\dim W }  \cdot Q^{}\left(   {f_{1}'\ot  f_{2}'\over f_{3}'\ot f_{4}' } \right).
\eeqq


By the multiplicity-one principle,  Theorem \ref{GZ thm} for special quadruples implies Theorem  \ref{GZ thm} in general, since under our assumptions the functional $Q$ is non-vanishing on special quadruples: this again follows from Proposition \ref{compare Qs} and Lemma \ref{6444}.1.(a)-(b). 
\end{proof}


\subsubsection{Comparison of $p$-adic $L$-functions}
We describe how, upon restricting  $ \calL_{p}(\cV^{\sharp})$ to the cyclotomic line through a point of trivial weight, we recover the $p$-adic $L$-function of \cite{dd-pyzz, nonsplit}. 
 


\begin{lemm} \lb{compare Lp} 
Let $z=(x, y)\in \cE_{\G_{0}}^{\ord, \rm cl}\ts \cE_{\H}^{\ord, \cl}$ be a point corresponding to a representation $\pi_{0 ,x}\otimes \chi_{y}$ of weights $(0; (2, \ldots, 2))$, $(0;0, \ldots, 0)$. 
 Let $A$ denote the modular abelian variety attached to $\pi_{0,x}$, and let  $$\calL_{p}( V_{(A , \chi)})\in \mathscr{K}(\cE_{\rm Z})$$
  be
  the $p$-adic $L$-function of \cite[Theorem A]{nonsplit}.
   Consider the map 
 \beqq
 j_{x, y}\colon \cE_{\rm Z}&\to \{x\} \ts \cE_{\H}^{\ord} \subset \cE_{\G_{0}}^{\ord, \rm cl}{\ts} \cE_{\H}^{\ord}\\
 \chi_{F} &\mapsto \pi_{0, x}\otimes \chi_{y}\cdot \chi_{F\circ N_{E_{\A}^{\ts}/F_{\A}^{\ts}}}.\eeqq
 Then 
\beq
\lb{LL'} \calL_{p}( V_{(\pi_{0} , \chi)}):=  j_{(x, y)}^{*}  \calL_{p}(\cV^{\sharp}) &=   \calL_{p}( V_{(A, \chi)}).\eeq  
\end{lemm} 
\begin{proof}  This is immediate from the respective  interpolation properties. (Note that the first equality in
  \eqref{LL'} is just a reminder of \eqref{def Lp}.)
\end{proof}






\subsubsection{Interpolation argument and proof of the main theorems}
Let $\X \subset \cE^{\ord}_{K^{p}}$ be a locally distinguished Hida family for $(\G\ts \H)'$, as in Definition \ref{hida loc dist}. Fix a level $K^{p}{}'\subset K^{p}$.  Let 
$$\X'=\X'_{K^{p}{}'}:=\X^{(6)}\cap \X^{(3,f)}\supset \X^{\cl}$$  be  the intersection of the open subsets of $\X$ of Theorem \ref{X6} and Corollary \ref{X3f}. 


Recall that we denote by $\X^{\cl, W}$ the set of classical points of weight $W$, and by $\X^{\cl, \textup{ n-exc}}$ the set of non-exceptional classical points. When $W=\Q_{p}$ is the trivial weight, we  also define
$$\X^{\cl,\ \textup{p-crys}, \  \Q_{p}, \textup{ram}} \subset \X^{\cl, \textup{p-crys}, \, \Q_{p}} \subset  \X^{\cl, \text{n-exc}} $$
by the following conditions on the classical point $z=(x, y)$ (equivalently, on the representation $\Pi_{z}$):
\begin{enumerate}
\item[(p-crys)] for all $v\vert p$, the representation $V_{x|G_{F_{v}}}$ is potentially crystalline (equivalently, $\pi_{x,v}$ is a principal series; the second inclusion above follows from Lemma \ref{6444});
\item[(ram)] $\chi_{y,p}$ is \emph{sufficiently ramified} in the following sense: let $r_{v}\geq1$ be minimal such that 
 $1+\vpi_{v}^{r_{v}} \OO_{F_{v}}$ is contained in the kernel of $\omega_{x,v}$, and let $U_{F,p}^{\circ}=\prod_{v\vert p}   (1+\vpi_{v}^{r_{v}} \OO_{F_{v}})$; then $\chi_{y, p}$ is is nontrivial on $N_{E_{p}/F_{p}}^{-1}(U_{F,p}^{\circ})\cap \OO_{E_{p}}^{\ts}$.
\end{enumerate}


 \begin{lemm} \lb{dense'}
 The subset $\X^{\cl,\ \textup{p-crys}, \  \Q_{p}, \textup{ram}}  \subset \X'$ is dense. 
 \end{lemm}
 \begin{proof} 
Denote by ${\rm p}_{\G} \colon \X' \to \cE_{\G}$ the natural projection.
 If  `$?$' is any relevant decoration, let  $\X_{\G}^{?}:= {\rm p}_{\G}(\X^{?})$; for  $x\in \X_{\G}^{\cl}$,  let $\X_{x, \H}^{?}:= {\rm p}_{\G}^{-1}(x)\cap  \X^{?}$.

 
 
For each $x\in \X_{\G}^{\cl, \textup{ p-crys, } \Q_{p}}$, the set 
$ \X_{x,\H}^{\cl,  \ \Q_{p},   \textup{ ram }}$ contains contains all but finitely many points in 
 $ \X_{x, \H}^{\cl, \, \Q_{p}}$, which is dense in  $\X_{x, \H}$. Thus the closure of $\X^{\cl,\ \text{p-crys}, \  \Q_{p}, \text{ ram}}$ contains all of $\X^{\cl,\ \text{p-crys},\  \Q_{p}}$.  
 
 Now we observe that $\X_{\G}^{\cl, \ \text{p-crys},\  \Q_{p}} = \mathscr{Y}_{\G}\cap \X_{\G}^{\cl,\  \Q_{p}}$ for the open subset  $$\mathscr{Y}_{\G}:=\X_{\G}- \{ x\ | \cV_{\G, v|x }^{+}(-1)  \cong \cV_{\G, v|x}^{-} \text{for some $v\vert p$}\} ,$$ which is non-empty as it contains $\X_{\G}^{\cl, W_{\G}}$ for any representation $W_{\G}$ whose partial weights are all $\geq 3$ (cf. the proof of Lemma \ref{6444}.1 (c)). 
Therefore  $\X^{\cl,\ \text{p-crys},\  \Q_{p}}$ is the intersection of the  non-empty open ${\rm p}_{\G}^{-1}(\Y_{\G})$ with  $\X^{\cl, \, \Q_{p}}$, which  is dense in $\X'$ by Lemma \ref{dense}.  We conclude that   $\X^{\cl,\ \text{p-crys},\  \Q_{p}}$ and  $\X^{\cl,\ \text{p-crys}, \  \Q_{p}, \text{ ram}}$ are dense in $\X'$. 
 \end{proof}
 
 % up to replacing it by $\X'\cap \iota(\X')$ we may assume it is $\iota $-stable. 







  



\begin{prop}
\source{universal-p-adic-gross-zagier-formula}
\notready \lb{final lemma}
The following are equivalent.
\begin{enumerate}
\item[1.] \lb{it1} Theorem \ref{ugz thm} holds over $\X'_{K^{p}{}'}$ for all $K^{p}{}'\subset K^{p}$.

\medskip



\item[2a.] \lb{it3a}
Theorem \ref{GZ thm'} holds for all representations $\Pi$ corresponding to points of $\X^{\cl}$ satisfying $\textup{(wt)}$.
\item[2b.] \lb{it3c}
Theorem \ref{GZ thm'} holds for all representations $\Pi$ corresponding to points of $\X^{\cl, \textup{ p-crys}, \ \Q_{p}, \textup{ ram}}$.

\medskip

\item[3a.] \lb{it4a} Theorem \ref{GZ thm}   holds for all representations $\Pi$ corresponding to points of $\X^{\cl, \textup{ n-exc}}$ satisfying $\textup{(wt)}$.
\item[3b.] \lb{it4c} Theorem \ref{GZ thm}  holds for all  representations $\Pi$ corresponding to points of 
$\X^{\cl, \textup{ p-crys}, \ \Q_{p}, \textup{ ram}}$.
\end{enumerate}
\end{prop}
\begin{proof}
 For any point   $z\in \X^{\cl, \textup{n-exc}}$ satisfying $\textup{(wt)}$,  denote  by $\Pi_{z}$ the associated automorphic representation, by $V_{z}$  the associated Galois representation . We have   proved the following specialisation-at-$z$ properties of objects defined  over (open subsets of)  $\X$:
\begin{itemize}
\item the $\OO_{\X'}$-module $\Pi_{H'_{\Sg}}^{K^{p}{}', \ord}$ (respectively $(\Pi_{H'_{\Sg}}^{K^{p}{}', \ord})^{\iota}$  specialises to $\Pi_{z,H'_{\Sg}}^{K^{p}{}', \ord}$ (respectively  $\Pi_{z,H'_{\Sg}}^{\vee, K^{p}{}', \ord}$), by Proposition \ref{cofinal}.\ref{Piord} (and the definition of the involution $\iota$);
\item the Galois representation $\cV$  (respectively $\cV^{\iota}$) and its ordinary filtrations  specialise to $V=V_{\Pi}$ (respectively $V_{\Pi^{\vee}}$)  with its ordinary filtrations, by construction (Proposition \ref{galois2});
\item there is a natural map $\wtil{H}^{1}_{f}(E, \cV)_{|z}\to \wtil{H}^{1}_{f}(E, V)$;
\item  the class $\cP_{K^{p}{}'}$ specialises to the restriction of $P^{\ord}_{\Pi_{z}}$ to $\Pi_{z,H'_{\Sg}}^{K^{p}{}', \ord}$ under the above map, by  Theorem   \ref{theo heeg univ} whose proof is completed in \S~\ref{645};
\item the product of local terms $\cQ$ specialises to the restriction of $Q^{\ord} $ to the spaces of $H'_{\Sg}$-coinvariants, $K^{p}{}'$-invariants in $\Pi_{z}^{\ord}$, $\Pi_{z}^{\vee,\ord}$, 
 by Theorem \ref{X6}.
\end{itemize}
Let us  complete the proof  that either side of Theorem \ref{ugz thm} specialises to $1/2$ times the corresponding side of Theorem \ref{GZ thm'}.  Consider the diagram $\X_{0}\to \X_{0}^{\sharp}\to \cE_{\rm Z}$. It is \emph{not} a product, even Zariski-locally; however  the
conormal sheaf is trivial. (This is dual to the fact that $\G\ts\H\to (\G\ts\H)'$ is a ${\rm Z}$-torsor for the \'etale topology but  {not} for the Zariski topology.)  The    immersion $\X\hat{\ts} \cE_{\rm Z}\to \X^{\sharp}$ given by `$(\Pi, \chi_{F})\mapsto \Pi\ot \chi_{F}$' induces the map on conormal sheaves 
$$\calN_{\X/\X^{\sharp}}^{*}=\OO_{\X}\hat{\ot} \Gamma_{F} \to
\calN_{\X/\X\ts \cE_{\rm Z}}^{*} = \OO_{\X}\hat{\ot} \Gamma_{F} $$
  that is multiplication by $1/2$
 under the canonical identifications. Hence:
\begin{itemize}
\item the $p$-adic height pairing ${h}_{\cV}=h_{\cV^{\sharp}|\X}$ specialises to 
${1\over 2} h_{V}={1\over 2} h_{V\ot \chi_{F, {\rm univ}|G_{E} }| z}$, by \S~\ref{sec: ht};
\item the derivative ${\rm d}^{\sharp} \calL_{p}(\cV)$   specialises to ${1\over 2} \calL_{p}'(V, 0)$ in $\Q_{p}(z)\hat{\ot} \Gamma_{F}$,  by the definition in \eqref{def Lp}.
\end{itemize}

We may now  complete the proof. By the specialisation properties summarised above, we have $1.\Rightarrow 2.a$ ($\Rightarrow 2.b$). By Lemma \ref{B Bord}, we have $2.a \Rightarrow 3.a$, $2.b \Leftrightarrow 3.b$.  
 By Lemma \ref{dense'} and the specialisation properties, $2.b\Rightarrow 1.$
\end{proof}


\begin{proof}[{Proof of Theorems \ref{GZ thm},  \ref{ugz thm}, and \ref{GZ thm'}}] The first assertion of Theorem \ref{GZ thm} was proved in Proposition \ref{in H1f}.
For a representation $\Pi$ of trivial weight satisfying the conditions (p-crys), (ram), the  formula of Theorem \ref{GZ thm} is \cite[Theorem B]{nonsplit} (cf. Lemma \ref{compare Lp}).
  By Proposition \ref{final lemma}, this implies  Theorem \ref{ugz thm} and the general case  of Theorems \ref{GZ thm},   \ref{GZ thm'}.  
\end{proof}



\subsubsection{Applications to non-vanishing / 1: self-dual CM families} \lb{sec: nvCM}
We prove the generic non-vanishing result of  Theorem \ref{nv CM}. Recall that $\Y$ is a component of the subvariety $\cE_{\H}^{\ord , \rm sd}\subset\cE_{\H}^{\ord}$ cut out by the condition $\chi_{F_{\A}^{\ts}}=\eta \chi_{{\rm cyc}, F}$, and such that $\eps(\chi_{y}, 1)=-1$ generically along $\Y$. 


\begin{proof}[Proof of Theorem \ref{nv CM}] Recall that a $p$-adic CM type of $E$ over  $\baar{\Q}_{p}$ is a choice   $\Sg$ of a place $w\vert v$ of $E$ for each place $v\vert p$ of $F$ (we identify primes above $p$ with embeddings into   $\baar{\Q}_{p}$). For each of the $p$-adic CM types $\Sg$ of $E$ and each connected component $\Y^{\sharp}$ of $\cE^{\ord}_{\H}$,  there is a Katz $p$-adic $L$-function
$$L_{\Sg} \in \OO(\Y^{\sharp}).$$
It is characterised (see \cite{katz, HT}) by its  values at the subset $\Y^{\sharp, \cl.\Sg}\subset \Y^{\sharp, \cl}$ of those $y$ such that  the algebraic part of $\chi_{y}$ is   $t\mapsto\prod_{\sg\in \Sg} \sg(t)^{w} \sg(t/t^{c})^{k_{\sg}}$ for integers $w,  k_{\sg}$ such that either $w\geq 1$, or  $w<1$ and $w+k_{\sg}>0$ for all $\sg\in \Sg$. The interpolation property relates $L_{\Sg}(y)$ to $L(1, \chi_{y})$. It is easy to see that for  a given $y\in \cE^{\ord, \rm sd}_{\H}$, there is a unique CM type $\Sg$ such that $y$ belongs to the interpolation subset of $L_{\Sg}$. For such $y$ and $\Sg$,  we denote by $L_{p}(\chi_{y}, s)\in \OO(\cE_{{\rm Z}/\Q_{p}(y)})$ the function $s \mapsto L_{\Sg}(y(s))$ where $\chi_{y(s)}=\chi\cdot \chi_{F,s}\circ N_{E/F})$. 


Consider now the setup of the theorem, and  let $\Y^{\sharp}\subset \cE_{\H}^{\ord}$ be the component containing $\Y$. 
By \cite{ashay}, under our assumptions the normal derivative $d^{\sharp} L_{\Sg}\in \mathscr{N}^{*}_{\Y/\Y^{\sharp}}$ is non-vanishing. Let $\Y'_{\Sg}$ be the complement of its zero locus, and let $$\Y':=\bigcap_{\Sg} \Y_{ \Sg}'.$$
 Let $y\in \Y^{\cl}\cap \Y'$, and let $\chi:=\chi_{y}$. It is easy to see that there is a unique $\Sg$ such that $y\in \Y^{\sharp, \cl , \Sg}$. 
 
 We \emph{claim} that there exists a a finite-order character $\chi_{0}\in \cE_{\H}^{\ord, \cl}$, such that  the character 
 $$\chi'= \chi^{c} \chi_{0}^{c}\chi_{0}^{-1} $$
 (that has  the same algebraic part as $\chi^{c}:=\chi\circ c$, and    defines a point $y'\in \cE_{\H}^{\ord, {\rm sd}, \cl}$) satisfies the following properties:
\begin{itemize}
 \item $L(1, \chi')\doteq L_{p}(\chi', 0)\neq 0$, where $\doteq $ denotes equality up to a nonzero constant;
 \item $H^{1}_{f}(E, \chi')=0$.
 \end{itemize}
 
 
 Granted the claim, we have a decomposition of $G_{E}$-representations
$$ (\chi\chi_{0}\oplus \chi^{c}\chi_{0}^{c})\ot \chi_{0}^{-1} =\chi \oplus \chi' $$ 
and a corresponding factorisation
$$ \calL_{p}(V_{(\pi_{0}, \chi_{0}^{-1})},s)\doteq L_{p}(\chi,s) L_{p}(\chi',s ), $$
 where $\pi_{0}=\theta(\chi\chi_{0})$ (the theta lift), and $\pi_{0}\otimes\chi_{0}^{-1}$ descends to a representation of $(\G_{0}\ts\H)'(\A)$.  It follows that  $\calL_{p}'(V_{(\pi_{0}, \chi_{0}^{-1})},0)\neq 0$. By Theorem \ref{BK thm}, we have a class 
 $$Z \in H^{1}_{f}({E}, V_{(\pi_{0}, \chi_{0}^{-1})}) \ot H^{1}_{f}({E}, V_{(\pi_{0}, \chi_{0}^{-1})}) =    (H^{1}_{f}({E},\chi) \otimes H^{1}_{f}({E},\chi^{c}))
 \oplus (H^{1}_{f}({E},\chi')\otimes H^{1}_{f}({E},\chi'^{c})) $$
  whose $p$-adic height is non-vanishing. Since $H^{1}_{f}({E},\chi')=H^{1}_{f}({E},\chi'^{c})) =0$, the class $Z$ is as desired. 


It remains to prove the claim. Let $\Y_{1}\subset \cE_{\H}^{\ord, {\rm sd}}$ be a component over which the anticyclotomic Main Conjecture is known -- that is, one containing a finite-order character satisfying the properties of \cite{hida-quad}). By applying  \cite[Lemma 2.5]{bur-d} to any character corresponding to a point of $\Y_{1}^{\cl}$, we find another component $\Y_{2}\subset \cE_{\H}^{\ord, {\rm sd}}$,  whose classical points correspond to characters $\chi_{2}$ with $$\eps(1, \chi_{2})=1;$$
 moreover from the proof in \emph{loc. cit.}  one sees that $\Y_{2}$ may be taken to still satisfy the assumptions of \cite{hida-quad}. Then the function $L_{\Sg c|\Y_{2}}$ is non-vanishing by \cite{hsieh-mu}; hence, by the density of classical points with a given weight, we may find $y'\in \Y_{2}^{\cl}$ corresponding to a character $\chi'$ satisfying the first among the required  conditions. By the anticylcotomic Main Conjecture for $\Y_{2}$ proved in \cite{hida-mc, hida-quad}, that is equivalent to the second condition. Finally, the ratio $\chi'/\chi^{c} $ is an anticyclotomic character (that is, trivial on $F_{\A}^{\ts}$), hence (\cite[Lemma 5.3.1]{28}) of the form $\chi_{0}^{c}\chi_{0}^{-1}$ for some finite-order character $\chi_{0}$.  
\end{proof}

\subsection{A universal  Waldspurger formula}\lb{sec: wald}
We describe  the complementary picture over locally distinguished families  attached to \emph{coherent} quaternion algebras over $F$.  Unexplained notions and  notation will be entirely parallel to what defined in the introduction. 

 Let $B$ be a totally definite quaternion algebra \emph{over $F$}, let  $\Sg$ be the set of finite places where $B$ is ramified, and let $\G_{/\Q}$ be the algebraic group with $\G(R)=(B\ot R)^{\ts}$ for any $\Q$-algebra $R$. Let $E$ be a CM quadratic extension of $F$, admitting an $F$-algebra embedding ${\rm e}\colon E\into B$  which we fix. We use the same symbols as in the introduction for the towers of Shimura varieties associated to the groups as in \eqref{list}.  Here, all those Shimura varieties  are $0$-dimensional.

 Let $L$ be a $p$-adic field and let $\Pi$ be an automorphic representation of $(\G\ts\H)'(\A):=(\G\ts\H)'(\A^{\infty})\ts (\G\ts\H')_{/\Q_{p}}$ over $L$ of weight $W$, by which we mean one occurring in $\H^{0}(\baar{Z}, \cW^{\vee})\ot W$. The normalised  fundamental class of $Y$ gives rise to an element $P\in \H_{0}(Z, \cW)^{\H'(\A)}$ and to an  $\H'(\A)$-invariant functional 
 $$P_{\Pi}\colon \Pi\to L, $$
 which may be nonzero only if $\Pi$ is locally distinguished by $\H'$. If $\Pi$ is ordinary, we again define $P^{\ord}:= P\tord$. 
 
 
 Let $\X$ be a locally distinguished  Hida family for $(\G\ts \H)'$, which via a Jacquet--Langlands  map is isomorphic to  a Hida family $\X_{0}$ for $(\G_{0}\ts\H)'$. For each compact open subgroup $K^{p}\subset (\G\ts\H)'(\A^{p})$, there is a `universal ordinary representation' $\Pi^{K^{p},\ord}_{H'_{\Sg}}$ of $(\G\ts\H)'(\A)$ over $\X$. As in Theorem \ref{theo heeg univ}, there exists an $\H'(\A^{p\infty})$-invariant, $\OO_{\X}$-linear functional 
\beq\lb{cP as functional2} \cP\colon  \Pi^{K^{p},\ord}_{H'_{\Sg}}\to \OO_{\X} \eeq
 interpolating (the restrictions of) $P_{\Pi_{z}}^{\ord}$ at all $z\in \X^{\cl}$ satisfying (wt). 
 
Starting from the natural pairings $\H_{0} (\baar{Z}_{K} , \cW)\ot \H_{0} (\baar{Z}_{K} ,\cW^{\vee})\to L$, we may define pairings $(\ , \ )_{\Pi}$ on each  representation $\Pi$ over a field  by the formula  \eqref{pair pi} (using the counting measure for ${\rm v}(K)$);   then we obtain modified pairings $(\ , \ )^{\ord}_{\Pi}$ on each $\Pi^{\ord}\ot \Pi^{\vee, \ord}$, and a  pairing $((\ ,  \ ))$ on the universal representations over $\X$, interpolating modified pairings  $(\ , \ )^{\ord}_{\Pi_{z}}$.

Finally, the functional $\cQ$, over an open $\X'\subset \X$ containing $\X^{\cl}$,  is also constructed as  in  \S~\ref{sec 44}; in the argument using the local Langlands correspondence, we use the rank-$2$ family of Galois representations pulled back from $\X_{0}$. 

 \begin{theoA}\label{wald univ} Let $\X$ be a locally distinguished Hida family for $(\G\ts\H)'$. 
\source{universal-p-adic-gross-zagier-formula}
 Abbreviate $\Pi^{(\iota)}:=  \Pi^{K^{p}{}', \ord, (\iota)}_{H'_{\Sg}}$, $\OO:=\OO_{\X}$, $\cK:=\cK_{\X}$. 

There is an
open subset $\X'\subset \X$ containing $\X^{\cl}$, such that 
$${\cP(f_{1})\cdot \cP^{\iota}(f_{2})    \over   ((f_{3}, f_{4}))    }
=
 \calL_{p}(\cV^{\sharp})_{|\X}
\cdot
\cQ\left( {  f_{1} \ot  f_{2}\over f_{3}\ot f_{4}}\right),
$$
 an equality of $\cK$-valued $\OO$-linear  functionals on $(\Pi\ot_{\OO} \Pi^{\iota})\ot_{\OO^{\times}} (\Pi\ot_{\OO} \Pi^{\iota})^{\ts, -1}$.
\end{theoA}
Similarly to \S~\ref{sec: final}, this universal formula follows from its specialisations at all classical points satisfying (wt); those are known by modifying the main result of  \cite{wald} as in Theorem \ref{GZ thm'}. 

 The formula essentially reduces the study of $\calL_{p}(\cV^{\sharp})_{|\X}$ to the study of the universal Waldspurger periods $\cP$. This is particularly interesting in the case of exceptional zeros, as we discuss next.

\subsection{Bertolini--Darmon conjectures and exceptional zeros} \lb{sec bd conj}
We first  formulate a conjecture on the behaviour of $\cP$ at a point $z\in \X^{\cl}$ and gather some old and new evidence in its favour.  In view of \S~\ref{ss 648}, the conjecture is often particularly interesting when $z$ is exceptional.  Then we deduce from our constructions and a known exceptional case of the conjecture a proof of Theorem \ref{nv exc}.










The conjecture requires some algebraic preliminaries.
 
 \subsubsection{Pfaffian regulators}  
 Let $L$ be a field  of characteristic $0$, and let $M$, $T$ be a finite-dimensional $L$-vector spaces.  Let $h\colon M\ot M\to T$ be a skew-symmetric pairing, and let $r=\dim_{L}M$.
 
 If $r$ is even, we define the Pfaffian regulator  $${\rm Pf}^{+}(M, h)\in ({\rm Sym}^{r/2}T )/ L^{\ts}$$
 to be the Pfaffian of the skew-symmetric matrix $h(x_{i}, x_{j})_{ij}$ for any $L$-basis $x_{i}$ of $M$. It is well-defined modulo $L^{\ts}$. 
 
 If $r$ is odd, we define an enhanced Pfaffian regulator 
  $${\rm Pf}^{-}(M, h)\in (M\ot {\rm Sym}^{(r-1)/2}T )/ L^{\ts}$$
as follows. It suffices to define $\partial^{e}{\rm Pf}^{-}(h) $ for any basis $\partial_{1}, \ldots, \partial_{d}$ of $T^{\vee}$ and all tuples $e=(e_{i})_{i=1}^{d}$ with $\sum_{i=1}^{d}e_{i}= (r-1)/2$; here $\partial^{e}:= \prod_{i=1}^{d} \partial_{i}^{e_{i}} $. Let $I\subset \{1, \ldots, d\}$ be the support of the tuple $e$. Let $M_{I}   \subset M$ be the sum, over $i \in I$,  of the radicals of the pairings $\partial_{i} h$. If $\dim M_{I}\geq 2$, we define $
\partial^{e} {\rm Pf}^{-}(h) :=0$. If $\dim M_{I}=1$ and $x\in M_{i}$ is a generator,  denote by $\baar{h}$ the induced pairing on $\baar{M}:= M/M_{I}$; we define 
$$\partial^{e} {\rm Pf}^{-}(M, h)  = x\ot \partial^{e} {\rm Pf}^{+}(\baar{M}, \baar{h}). $$

\begin{rema}
\source{universal-p-adic-gross-zagier-formula}
\notready\lb{rel reg} In the even case, we have of course ${\rm Pf}^{+}(M, h) ^{2}  = R(M, h)\in ( {\rm Sym}^{r}T )/ L^{\ts,2}$, where $R(M, h)$ is the discriminant of the pairing $h$. In the odd case, assume further given a symmetric bilinear pairing $h^{\sharp}\colon M\ot M\to T^{\sharp}$. Let $h'=h\oplus h^{\sharp}\colon M\ot M\to T':-T\oplus T^{\sharp}$, and let $R(M, h')\in {\rm Sym}^{r}( T')/L^{\ts, 2}$ be its discriminant. Then it is easy to verify that 
$$h^{\sharp}( {\rm Pf}^{-}(M, h), {\rm Pf}^{-}(M, h))\in (T^{\sharp}\ot {\rm Sym}^{r-1}T/) L^{\ts, 2}$$
 is the image of $R(M, h')$ under the natural projection.
\end{rema}

\begin{rema}
\source{universal-p-adic-gross-zagier-formula}
\notready\lb{ambi int} If $L$ is the fraction field of a domain $\OO$ and $M$ is endowed with an $\OO$-lattice, it is possible to lift the ambiguity in the definitions to an element of $\OO^{\ts}$. 
\end{rema}







 
 
 


\subsubsection{A conjecture \`a la Bertolini--Darmon} \lb{sec 733}
Let  $\G$ be either as in the introduction or as in \S~\ref{sec: wald}. We define a sign $\epsilon := -1$ in the former case and $\epsilon := +1$ in the latter case. Let $\X$ be a locally distinguished Hida family for $(\G\ts\H)'$, and let $$z\in \X^{\cl}$$ be a classical point.  We denote by $\cP$ the universal Heegner class (if $\epsilon=-1$) or toric period (if $\epsilon =+1$), viewed as a family of functionals as in \eqref{cP as functional}, \eqref{cP as functional2}, parametrised by a subset $\X'\subset \X$.  Assume that $\X'$ can be taken to be a neighbourhood  of $z\in \X$.  



Let  $T_{z}^{*}\X ={\frakm}_{z}/\frakm_{z}^{2}$ be the cotangent space to $\X$ at $z$ (where $\frakm_{z}\subset \OO_{\X, z}$ is the  maximal ideal), and for any $r\in \N$ let  
$$({\rm d}_{z}^{\parallel})^{r} \colon {\frakm}_{z}^{r}\subset \OO_{\X, z}\to \frakm_{z}^{r}/\frakm_{z}^{r+1} = {\rm Sym}^{r} T_{z}^{*}\X  $$
be the   natural projection. It is easy to see that the involution $\iota$ of \S~\ref{ssec invo} satisfies
${\rm d}_{z}^{\parallel}\iota =\id$ (whereas ${\rm d}_{z}^{\sharp} \iota =-\id$).  

Let $V=\cV_{|z}$, and let
 ${\rm c} \colon  \wtil{H}^{1}_{f}(E, V^{\iota}) \to  \wtil{H}^{1}_{f}(E, V)$ be the isomorphism induced by the adjoint action, on $G_{E, S}$, of a lift of the complex conjugation in $\Gal(E/F)$. Let 
\beq\lb{hpar}h^{\parallel}:= h_{z/\X}^{\Box} \colon \wtil{H}^{1}_{f}(E, V) \ot  \wtil{H}^{1}_{f}(E, V)  \to  T_{z}^{*}\X\eeq
 be the \nek--Venerucci height pairing as in \eqref{skewsym}.  Since the pairing of Proposition \ref{duality} is skew-hermitian, by Proposition \ref{prop ht2} the pairing $h^{\parallel}$ is skew-symmetric.
Define  $${\rm Pf}^{\, \parallel, \pm}(V):= {\rm Pf}^{\pm} (\wtil{H}^{1}_{f}(E, V), h^{\parallel}).$$



  \begin{customconj}
\source{universal-p-adic-gross-zagier-formula}
\notready{Pf} \lb{pf conj}  Let $\wtil{r}:= \dim_{L}\wtil{H}^{1}_{f}(E, V)$.
We have $(-1)^{\wtil{r}}=\epsilon$, the universal element  $\cP$ vanishes to order at least $\lfloor{\wtil{r}/2}\rfloor$ at $z$, and for any generator  $\wp\in (\Pi^{\ord})^{*, \H'(\A^{p\infty)}}$ and all $f\in \Pi^{\ord}$, we have
\beq\lb{pf conj eq}
({{\rm d}}^{\parallel}_{z})^{ \lfloor{\wtil{r}/2}\rfloor} \, \cP(f) = {\rm Pf}^{\, \parallel, \epsilon}(V) \cd \wp(f)\eeq
in  $[\wtil{H}^{1}_{f}(E, V)^{(1+\epsilon)/2}\ot_{\Q_{p}(z)} {\rm Sym}_{\Q_{p}(z)}^{{\lfloor{\wtil{r}/2}\rfloor}}T_{z}^{*}\X ]\, / \Q_{p}(z)^{\ts}$.
\end{customconj}


\subsubsection{Relation to the original conjectures of Bertolini--Darmon}\lb{ssec bd} Define $\X^{\rm a}:= \X\cap ( \{x\}\ts \cE_{\H'})$ (the anticyclotomic family), $\X^{\rm wt} := \X\cap (\cE_{\G}\ts\{y\})$ (the weight family). In the classical case when $E$ is imaginary quadratic, $\pi$ is associated with an elliptic curve over $\Q$, $\chi=\one$,  and we restrict to the anticyclotomic variable $\X^{\rm a}$, Conjecture \ref{pf conj} is a variant of conjectures of Bertolini--Darmon, surveyed in \cite{bdsurvey} and (in a form slightly closer to the one of the present work) in \cite[\S~4]{dd-exc}. The Bertolini--Darmon conjectures were partly generalised to higher-weight modular forms  in \cite{LV}.


\begin{rema}
\source{universal-p-adic-gross-zagier-formula}
\notready By using natural  $G_{E, S}$-stable lattices in $V$, or better lattices in $\wtil{H}^{1}_{f}(E, V)$ spanned by motivic elements, it is possible to define the Pfaffian regulators up to an ambiguity that is a unit in the ring of integers of a local field or of a number field (recall Remark \ref{ambi int}). It should then possible to refine the conjecture up to such ambiguity (by including the appropriate constants), as in the original works of Bertolini--Darmon (see \cite[Conjecture 4.2.1]{dd-exc}).

In view of Remark \ref{rel reg}, inserting Conjecture \ref{pf conj} into Theorem \ref{ugz thm} would yield a multivariable formula relating higher partial derivatives of $p$-adic $L$-functions  with suitable height regulators, in the spirit of the  Birch and Swinnerton-Dyer conjecture; the argument is the same as that of  \cite[Proposition 5.1.1]{dd-exc}. We plan to return to formulate such conjectural formulas in the appropriate generality in future work. 
\end{rema}





\subsubsection{Evidence for Conjecture \ref{pf conj} in low rank} The preliminary parity conjecture 
\beq \lb{par1} 
(-1)^{\wtil{r}}=\epsilon\eeq 
is known in many cases as a consequence of the work of \nek\ (see \cite[Theorem 12.2.3]{nek-selmer}. Indeed, the statement proved  in \emph{loc. cit.}, in view of the functional equation of $\calL(V_{(\pi, \chi)},s)$, is that 
\beq \lb{par2}
(-1)^{r}= \eps
\eeq
 where $r=\dim H^{1}_{f}(E, V)$ and $\eps =\eps (V)$. Now if $r^{\rm exc}= r^{\rm exc, \, s}+ {r}^{\rm exc,\, ns}$ denotes the number of exceptional  primes of $F$ above $p$ (respectively, the number of those that moreover are split or nonsplit in $E$),  by Lemma \ref{6444}.2 we have $\epsilon=\eps\cd (-1)^{r^{\rm exc,\, ns}}$, and $\wtil{r}= r+2r^{\rm exc,\, s}+ r^{\rm exc, \, ns}$. Thus \eqref{par2} is equivalent to \eqref{par1}.

Let us now review the conjectural vanishing and leading-term formula. Most of the available evidence is concentrated in the case where $V$ arises from the classical context of \S~\ref{ssec bd}, to which we restrict unless otherwise noted for the rest of this discussion. (All the results mentioned below hold under various additional assumptions, which we will not recall.)

 If $\X$ is replaced with $\X^{\rm a}$ (the original Bertolini--Darmon case), the conjecture is  known if $\calL(V_{(\pi, \chi)},s)$ vanishes  to order $0$ or $1$ at $s=0$, see \cite[Theorem 4.2.5]{dd-exc} and references therein, as well as \cite{LP}. Some of those results have been generalised to higher weight (\cite{IoS, Mas}) or to totally real fields (\cite{Hu, BGe, MB}).
   
 
  When $\X$ is replaced by $\X^{\rm wt}$, $p$ is inert in $E$, $V$ is exceptional, and $\epsilon =+1$, Bertolini--Darmon \cite{bd-hida} proved a formula  for  ${\rm d}_{z}^{\parallel} \cP_{|\X^{\rm wt}}$, which implies the projection of \eqref{pf conj eq} to $T_{z}^{*}\X^{\rm wt}$ when ${\rm ord}_{s=0}\calL(V_{(\pi, \chi)}, s)=1$. The interpretation of the formula of Bertolini--Darmon in terms of height pairings was observed by Venerucci (see \cite[Theorem 2.1 and Theorem 4.2.2]{ven}), whose work was a second important influence in the formulation of Conjecture \ref{pf conj}.  The Bertolini--Darmon formula was generalised to higher-weight modular forms by Seveso \cite{Sev}  and to elliptic curves over totally real fields by Mok \cite{Mok}.
  
  
  \subsubsection{Evidence  for Conjecture \ref{pf conj} in higher rank}
Lower bounds for the order of vanishing of $\cP_{|\X^{\rm a}}$ have been obtained in two recent works for $\epsilon =1$. In the context of elliptic curves over totally real fields, \cite[Theorem 5.5]{BGe} gives a  bound (that is coarser than predicted by Conjecture \ref{pf conj}) in terms of the number of exceptional primes. In a classical context (and if if $p$ splits in $E$),  Agboola--Castella \cite[Corollary 6.5]{Agb-Ca} prove a bound that is finer than that of Conjecture \ref{pf conj}. 
 (That refined  bound is predicted by Bertolini--Darmon;  it is related to some trivial  degeneracies of the anticyclotomic height pairing, as touched upon also in the paragraphs preceding \eqref{exppl} below.)

  
 
 Regarding  the formula of Conjecture  \ref{pf conj}, an interesting anticyclotomic case can be deduced from  the recent work \cite{FG}, as we now explain. 

 \paragraph{The work of Fornea--Gehrmann}
Suppose that $A$ is a modular elliptic curve over the totally real field $F$,  such that the set $S_{p}^{\rm exc}$
of places  $v|p$ where $A$ has multiplicative reduction consists of exactly  $r$ primes $v_{1}, \ldots, v_{r}$ inert in $E$.  %Assume also that the abelian variety $\Res_{F/\Q}A$ is simple.\footnote{This corresponds to the hypothesis of primitivity of $L(A,s)$ in \cite{FG, FGM}.} 
Let $\phi\colon\bigotimes_{i=}^{r} E_{v_{i}}^{\ts}\to \bigotimes_{i=1}^{r} {A}(E_{v_{i}})$  be the product of Tate unifomisations, and let  $\hat{\phi}$ be the induced map on $p$-adic completions.  
 
 
 One of the main constructions of \cite{FG} produces an explicit  element $Q_{A}\in \bigotimes_{i=1}^{r} E_{v_{i}}^{\ts}\hat{\ot} \Q_{p}$, which carries a precise conjectural relation to the arithmetic of $A$. Partition $S_{p}^{\exc}= S_{p}^{\exc, +}\sqcup S_{p}^{\exc, -}$ according to whether the multiplicative reduction is,  respectively, split or nonsplit, and let $r^{+}:= |S_{p}^{\exc, +}|$.  Denote   $ \hat{A}(E_{v})=A(E_{v})\hat{\ot}\Q_{p}$ and let $\hat{A}(E_{v})^{\pm}$ be its $\pm$-eigenspaces for the conjugation $c_{v}\in \Gal(E_{v}/F_{v})$. 
 
 
 
 Assume that $A_{E}$ has rank  $r_{A}\geq r$ and let $r_{A}^{+}$ be the rank of $A$. 
According to \cite[Conjectures 1.3, 1.5]{FGM}, we have\footnote{In \emph{locc. citt.}, some restrictive assumptions are made (in particular that $E$ is not CM), but the conjectures make sense even without those and indeed closely related conjectures appear in \cite{FG} without those assumptions. Moreover, our statement  slightly differs from the ones of \cite{FG}, which instead of postulating that $(x_{1}, \ldots, x_{r})$ is a basis, postulates that $\hat{\phi}(Q_{A})\neq 0 $ under the extra assumption that $\Res_{F/\Q}A$ is simple (equivalently $L(A, s)$ is primitive). Our slight reformulation appears more uniform and still addresses \cite[Remark 1.1]{FG} (cf. the comment following \cite[Conjecture 1.5]{FG}). }
 \beq \lb{conj FG}
  r_{A}>r \text{ or }  r_{A}^{+}+r^{+}\neq r &\Longrightarrow \hat\phi(Q_{A}) = 0, \\
  r_{A}= r \text{ and }  r_{A}^{+}+r^{+}=r &\Longrightarrow
\hat\phi(Q_{A}) \doteq \det( (\hat{x}^{\rm a}_{i,v_{j}})_{1\leq i, j\leq r}) \eeq
where $(x_{1}, \ldots, x_{r})$ is a basis of $A(E)_{\Q}$, and for $v\in S_{p}^{\exc , \pm} $ 
we denote by $\hat{x}_{v}^{\rm a}$ the eigen-projection  of $x\in A(E)$ to $ \hat{A}(E_{v})^{\mp}$. The symbol $\doteq $ denotes equality up to a constant in $\Q^{\ts}$. 

For  $V=V_{p}A_{E}$, assuming the finiteness of $\Sha(A_{E})[p^{\infty}]$ we have $\wtil{r}:=\dim  \wtil{H}^{1}_{f}(E, V)  =r+r_{A}$ (see Lemma \ref{6444}.1(a) or \eqref{exppl} below), and   the parity conjecture is known. Hence if $r_{A}\equiv r\pmod 2$,  which we henceforth assume,
  $V$ corresponds to a point $z$ of a locally  distinguished Hida family $\X$ of sign $\epsilon =+1$ (for a unique choice of the coherent quaternionic group $\G$).
 Let 
 $$\Gamma^{\rm a}:= (E_{\A^{\infty}}^{\ts}/F_{\A^{\infty}}^{\ts}\widehat{\OO}_{E}^{p, \ts}) \hat{\ot}\Q_{p}\cong T_{z}^{*}\X^{\rm a},  
 \qquad  \ell^{\rm a}=\prod_{w} \ell^{\rm a}_{w}\colon E_{\A^{\infty}}^{\ts}/E^{\ts} \to \Gamma^{\rm a}$$
  be the natural projection, and let 
  $\ell_{\rm exc, \ot}^{\rm a}:=\bigotimes_{i=1}^{r}\ell^{\rm a}_{v_{i}}\colon \bigotimes_{i=1}^{r} E_{v_{i}}^{\ts}\hat{\ot}\Q_{p} \to {\rm Sym}^{r}\Gamma^{\rm a}$.  Denoting by ${\rm d}_{z}^{\rm a}$ the component of ${\rm d}_{z}^{\parallel}$ along $\X^{\rm a}$,  
Theorem A of \cite{FG}  (which relies on the aforementioned lower bound from \cite{BGe}) proves 
\beq \lb{eq FG} ({\rm d}_{z}^{\rm a})^{ r} \cP(f^{\circ}) \doteq \ell_{\rm exc, \ot}^{\rm a}(Q_{A}) \eeq
for a suitable test vector $f^{\circ}\in \Pi^{\ord}$. 

 \paragraph{Comparison with Conjecture \ref{pf conj}}
We show that granted \eqref{conj FG} and the finiteness of $\Sha(A_{E})[p^{\infty}]$, the formula \eqref{eq FG} is equivalent to the conjectured \eqref{pf conj eq}. 







Let $h^{\rm a}\colon \wtil{H}^{1}_{f}(E, V) \ot \wtil{H}^{1}_{f}(E, V) \to T_{z}^{*}\X^{\rm a}=\Gamma^{\rm a}$
be the projection of $h^{\parallel}$, and let 
$${\rm Pf}^{\, \rm a, +}(V) ={\rm Pf}^{+} ( \wtil{H}^{1}_{f}(E, V) , h^{\rm a}) \in {\rm Sym}^{r} \Gamma^{\rm a}.$$
Let $\wtil{H}^{1}_{f}(E, V)^{\pm}$ be the eigenspaces for the complex conjugation $c\in \Gal(E/F)$.   Since $h_{z/\X^{\rm a}}$ enjoys the $c$-equivariance property $h_{z/\X^{\rm a}}(cx, cx')=c. h_{z/\X^{\rm a}}(x, x')$ and $c$ acts by $-1$ on $\Gamma^{\rm a}$, by construction the pairing $h^{\rm a}$ satisfies $h(cx, cx')= -h(x, x')$,  so that each of  $\wtil{H}^{1}_{f}(E, V)^{\pm}$ is $h^{\rm a}$-isotropic. In particular, Conjecture \ref{pf conj} agrees with \eqref{eq FG} and \eqref{conj FG} that, in the first case of the latter, we have $({\rm d}_{z}^{\rm a})^{r}\cP=0$.

Assume now we are in the second case of \eqref{conj FG}, so that each of  $\wtil{H}^{1}_{f}(E, V)^{\pm}$  
has dimension $r$ and $h^{\rm a}$ need not be degenerate.
For $1\leq i\leq r$,  let $q_{i}=q_{v_{i}}\in E_{v_{i}}^{\ts}$ be a Tate parameter for $A/E_{v}$. By
  \cite[\S~7.14]{nekheights}
 we explicitly have
\beq \lb{exppl}
\wtil{H}^{1}_{f}(E, V) &=  H^{1}_{f}(E, V) \oplus \bigoplus_{i=1}^{r} \Q_{p}\cd[q_{v_{i}}] \\
 h^{\rm a}(q_{v}, q_{v'})=0,\qquad  h^{\rm a}(x, q_{v}) &= \log^{\rm a}_{A, v} (\hat{x}_{v}) , \qquad x\in A(E),\quad  v\neq v',
\eeq
where $\log_{A, v}^{\rm a}\colon \hat{A}(E_{v})\ot \Q_{p} 
\cong E_{v}^{\ts}/q_{v}^{\Z}\hat{\ot}\Q_{p} 
\cong  \OO_{E_{v}}^{\ts}\hat{\ot}\Q_{p} \stackrel{\ell_{v}^{\rm a}}{\longrightarrow} \Gamma^{\rm a}$. Note that $\log_{A, v}^{\rm a}$ factors through $x\mapsto \hat{x}_{v}^{\rm a}$. 

Up to changing the basis  $x_{i}$ of $A(E)_{\Q}$ and reordering the $v_{i}$,  we may assume that the basis 
 $$x_{r^{+}+1}, \ldots, x_{r}, q_{1}, \ldots, q_{r^{+}}, \qquad  x_{1}, \ldots, x_{r^{+}}, q_{r^{+}+1}, \ldots, q_{r} $$
of $\wtil{H}^{1}_{f}(E, V)$ is the concatenation of a basis of  $\wtil{H}^{1}_{f}(E, V)^{+}$ and a basis of $\wtil{H}^{1}_{f}(E, V)^{-}$, respectively. Using this basis, \eqref{exppl},  and the identity ${\rm pf} \smalltwomat {}{M}{-M^{\rm t}}{} =\pm \det M$, we have\footnote{All the equalities to  follow ignore signs and in fact, by our coarse definitions, only make sense at best  up to $\Q^{\ts}$.}
$${\rm Pf}^{\, \rm  a, +}(V)=\det M= \det M_{1}\det M_{2},$$
 where the $r\ts r$ matrix $M$ is block-left-upper-triangular with anti-diagonal blocks $M_{k} = (\log_{A,v_{j}}^{\rm a}(x_{i}))_{i, j \in I_{k}}$ for $I_{1}=\{ r^{+}+1,\ldots, r\}$, $I_{2}=\{ 1, \ldots, r^{+}\}$. 


On the other hand, we note that under \eqref{conj FG}, we have $\ell_{\rm exc, \ot}^{\rm a}(Q_{A}) = \det  N$ where   the $r\ts r$ matrix $N_{ij}= \log^{\rm a}_{A,v_{j}}(\hat{x}_{i, v}) = h^{\rm a}(x_{i}, q_{j})$ is block-diagonal with  blocks $N_{1}=M_{2}$, $N_{2}=M_{1}$. Thus $\ell_{\rm exc, \ot}^{\rm a}(Q_{A})  = {\rm Pf}^{\, \rm  a, +}(V)$, and \eqref{eq FG} is equivalent to \eqref{pf conj eq}.










\subsubsection{Applications to non-vanishing / 2: exceptional families} \lb{sec G} We prove Theorem \ref{nv exc}.\footnote{A less interesting variant of it was sketched in \cite{nonsplit}.} Recall that $\X_{0}$ is a Hida family for ${\rm PGL}_{2/\Q}$, that contains a classical point $x_{0}\in \X(\Q_{p})$ corresponding to an elliptic curve $A$ with split multiplicative reduction at $p$ satisfying $L(A, 1)=L(V_{p}A, 0)\neq 0$. 

\begin{proof}[Proof of Theorem \ref{nv exc}]
 Let $E$ be an imaginary quadratic field, with associated quadratic character $\eta$, satisfying the following:  $p$ is inert in $A$, all other primes dividing the conductor of $A$ split in $E$, and the twisted $L$-value $L(A, \eta, 1)\neq 0$. Then $A$ has split multiplicative reduction over $E$ with Tate parameter
  $$q=q_{A}\in E_{p}^{\ts}.$$ 
Let $\Omega_{A_{E}}\in \bC^{\ts}$ be the N\'eron period, and let 
  let $\H:={\rm Res}_{E/\Q}{\bf G}_{m}$. 

By construction,  $\eps_{v}(V_{p}A_{E})=1$ for all finite $v\nmid p$, hence the Hida family $\X\subset \cE_{(\GL_{2}\ts\H)'}^{\ord}$ containing the image of $\X_{0}\ts \{\one\}$ is locally distinguished. Let $\X^{\sharp}\subset  \cE_{(\GL_{2}\ts\H)}^{\ord}$ be the Hida family containing $\X$.
Let $\Pi$ be the universal ordinary representation over $\X$ and let $f\in \Pi$ be such that $f_{|x_{0}}$ is a test vector (that is, a vector not annihilated by any $\H'(\A^{p})$-invariant  functional $\lm\colon \Pi_{|x_{0}}\to \Q_{p}$). Let $\cP_{0, E}$ be the pullback of $\cP(f)$ to $\X_{0}$, and let $\cP_{0}:={1\over 2}\Tr_{E/\Q}\cP_{0, E}$.  By Corollary \ref{X3f}, 
$$\cP_{0,E}\in \wtil{H}_{f}^{1}(G_{E}, \cV_{0}),\qquad \cP_{0}\in  \wtil{H}_{f}^{1}(G_{\Q}, \cV_{0}).$$



By the main result of \cite{BDrigid} (as reformulated in \cite[Theorem 5.4, \S~5.2]{bdsurvey}),\footnote{In the works of Bertolini--Darmon, an explicit test vector $f$ is chosen;  cf. \cite{dd-exc} for more details on bridging the setups.} there is a constant $c\in \Q_{p}^{\ts}$ such that 
$$\qquad \cP_{0, E}(x_{0})\ot \cP_{0, E}^{\iota}(x_{0})=c\cdot {L(A_{E}, 1)\over \Omega_{A_{E}} }\cdot [q] \ot [q]\qquad \text{in }  \wtil{H}^{1}_{f}({\Q}, V_{p}A_{E})\ot  \wtil{H}^{1}_{f}({\Q},V_{p} A_{E})  $$
using the  description
$ \wtil{H}^{1}_{f}({\Q}, V_{p}A_{E})= \Q_{p}\cdot [q] \oplus H^{1}_{f}({\Q}, V_{p}A_{E})$ as in \eqref{exppl}.

In particular, $\cP_{0, E}(x_{0})= \cP_{0}(x_{0})$ is a nonzero multiple of $[q]$, which   is $\Gal(E/\Q)$-invariant. Hence $\cP_{0, E}$ and $\cP_{0}$ are  non-vanishing.  Then by  \cite{nek-koly},  $\wtil{H}^{1}_{f}({\Q}, \cV_{0})$ has generic rank~1. 


Moreover, 
\beq\lb{GSa}
h_{\cV_{0}/\cV_{0}^{\sharp}}(\cP_{0} , \cP_{0}^{\iota})(x_{0}) =c\cdot {L(A_{E}, 1)\over \Omega_{A_{E}} }\cdot h([q], [q]) = c\cdot \ell (q) \cdot {L(A_{E}, 1) \over \Omega_{A_{E}}}\in \Gamma_{\Q}\ot_{\Z_{p}}\Q_{p},\eeq
where $\ell\colon \Q_{p}^{\ts}\to \Gamma_{\Q}$ is the universal logarithm (see again  \cite[\S~7.14]{nekheights} for the  second equality). By \cite{st et}, the right-hand side is nonzero, hence $h_{\cV_{0}/\cV_{0}^{\sharp}}(\cP_{0}, \cP_{0}^{\iota})\neq 0$. 
\end{proof}
\begin{rema}
\source{universal-p-adic-gross-zagier-formula}
\notready \lb{rmk GS}
As noted in \cite{dd-exc, nonsplit}, the combination of Theorem \ref{ugz thm} (or rather Theorem \ref{GZ thm'}) and a precise form of \eqref{GSa} gives a new proof ot the following  theorem of Greenberg--Stevens \cite{GS}: for $A_{/\Q_{p}}$ an elliptic curve of split multiplicative reduction at $p$ and $L_{p}(V_{p}A)\in \Z_{p}\llb \Gamma_{\Q}\rrb_{\Q_{p}}$ its $p$-adic $L$-function,    $$L_{p}'(V_{p}A, 0)= {\ell(q)\over{\rm ord}_{p}(q)} \cdot {L(A, 1)\over \Omega_{A}}.$$
\end{rema}

















  
\appendix
